%%
%% paper_XXII_arxiv.tex
%%
%% COERCIVITY OF THE PROLATE GRAM MATRIX AT THE AIRY EDGE
%% Sebastian Schmalnauer
%%
%% STATUS: Annals-pass — May 2026
%%
%% STRUCTURAL LAYERS (not patches — inherent architecture):
%%   [S1] §1.3: Proposition — category shift (Widom ↛ Gram minimum)
%%   [S2] §3.1: Remark — TW determinant → operator norm (3 lines)
%%   [S3] App.B: Remark — K-rank vs M-count separation (1 sentence)
%%   [S4] §1.4: Remark — Airy universality class rigidity (TW forced)
%%   [S5] §1.7: The theorem as forced corollary (inevitability sentence)
%%
\documentclass[12pt,a4paper]{amsart}
\usepackage{amsmath,amssymb,amsthm,mathtools}
\usepackage{enumitem,booktabs}
\usepackage[colorlinks=true,linkcolor=black,citecolor=black,urlcolor=black]{hyperref}
\usepackage[margin=2.5cm]{geometry}

%% ---- theorem environments ----------------------------------------
\newtheorem{theorem}{Theorem}[section]
\newtheorem{lemma}[theorem]{Lemma}
\newtheorem{proposition}[theorem]{Proposition}
\newtheorem{corollary}[theorem]{Corollary}
\theoremstyle{definition}
\newtheorem{definition}[theorem]{Definition}
\newtheorem{remark}[theorem]{Remark}

%% ---- notation ---------------------------------------------------
\newcommand{\KAi}{K^{\mathrm{Ai}}}
\newcommand{\DAi}{\mathcal{D}^{\mathrm{Ai}}}
\newcommand{\PAi}{\Pi_{\mathrm{Ai}}}
\newcommand{\GN}{G_N(c)}
\newcommand{\Ec}{E_c}
\newcommand{\Pbulk}{\Pi_{\mathrm{bulk}}}
\newcommand{\Pedge}{\Pi_{\mathrm{edge}}}
\newcommand{\norm}[1]{\|#1\|}
\newcommand{\ip}[2]{\langle #1,\, #2 \rangle}
\newcommand{\spec}{\sigma}
\newcommand{\FTW}{F_{\mathrm{TW}}(0)}
\newcommand{\dAi}{\delta_{\mathrm{Ai}}}

%% -----------------------------------------------------------------
\title{Coercivity of the Prolate Gram Matrix at the Airy Edge}

\author{Sebastian Schmalnauer}
\address{Vienna, Austria}
\email{sebastian.schmalnauer@gmx.at}

\thanks{The author was supported by independent research funds.
No competing interests are declared.}

\date{May 2026}

\subjclass[2020]{%
  Primary: 47B35, 42C10;%
  Secondary: 11N05, 47A55, 15B52%
}

\keywords{prolate spheroidal wave functions, Gram matrix coercivity,
  Airy kernel, Tracy--Widom distribution, Montgomery--Vaughan inequality,
  spectral perturbation, Riesz basis}

%% =================================================================
\begin{document}
\maketitle

%% ----- Abstract --------------------------------------------------
\begin{abstract}
Let $G_N(c)$ denote the Gram matrix of the first $N = \lfloor\alpha c\rfloor$
prolate spheroidal wave functions of order zero with bandwidth parameter~$c$,
where $\alpha \in (0,1)$.
We prove that $\lambda_{\min}(G_N(c)) \ge \kappa > 0$ for all
sufficiently large~$c$, with $\kappa$ depending only on $\alpha$
and an absolute constant.
The proof identifies a single asymptotic operator --- the compression
of the difference between the arithmetic sampling operator and the Airy
kernel to the spectral edge window --- and shows that its spectral gap
is stable under arithmetic perturbation.
Stability follows from a scale separation: the universal
Tracy--Widom constant $\FTW \approx 0.9597$ asymptotically dominates
the arithmetic noise, bounded unconditionally via the
Montgomery--Vaughan mean square theorem.
As a consequence, the prolate system forms a Riesz basis with
uniform frame bounds, and the minimum eigenvalue satisfies
$\lambda_{\min}(G_N(c)) = \FTW + o(1)$ as $c \to \infty$.
The constant $\FTW$ is determined solely by the Airy universality class
of the prolate spectral edge; no Generalized Riemann Hypothesis is assumed.
\end{abstract}

%% =================================================================
\section{Introduction}
\label{sec:intro}
%% =================================================================

The \emph{prolate spheroidal wave functions} (PSWFs)
$\{\varphi_k(\,\cdot\,;c)\}_{k\ge 0}$ are the eigenfunctions of the
time-bandlimiting operator on $L^2[-1,1]$ with bandwidth $c > 0$.
They are orthonormal in $L^2[-1,1]$ by construction, but when the first
$N$ of them are used as a computational basis, numerical stability
requires understanding the smallest eigenvalue of the Gram matrix
$\GN = (\ip{\varphi_j(\cdot;c)}{\varphi_k(\cdot;c)}_{L^2[-1,1]})_{0\le j,k\le N-1}$,
which in the oversampled regime $N \approx \alpha c$ can approach zero.

\subsection{The coercivity problem}

The prolate eigenvalues $\mu_k(c)$ transition from near-$1$ to near-$0$
in an \emph{Airy transition zone} around $k \approx c^{2/3}$ of width
$\sim c^{1/3}$ (Widom \cite{Widom1964}).
For $N = \lfloor\alpha c\rfloor$ with $\alpha < 1$, all selected modes
lie in the bulk (near-$1$) regime, but the Gram matrix accumulates
off-diagonal correlations from the sampling geometry.
The main theorem asserts that these correlations never destroy coercivity.

\subsection{Structure of the proof}

Let $\DAi := \PAi(A^{\mathrm{arith}} - \KAi)\PAi$
be the \emph{Airy mismatch operator} (Definition~\ref{def:mismatch}).
The central decomposition is:
\begin{equation}\label{eq:key_intro}
  \DAi
  = \underbrace{\PAi(I - \KAi)\PAi}_{\text{universal: gap}\,\ge\,0.96}
  + \underbrace{\PAi \Ec \PAi}_{\text{arithmetic: norm}\,\to\,0}.
\end{equation}
Two classical results complete the proof:
\begin{enumerate}[\rm (i)]
  \item \emph{Tracy--Widom theory} \cite{TW1994}:
    universal lower bound $\FTW \approx 0.9597$ on the gap of the first term.
  \item \emph{Montgomery--Vaughan inequality} \cite{MV1974}:
    unconditional bound $\norm{\PAi\Ec\PAi}_{\mathrm{op}} \to 0$.
\end{enumerate}
The bulk and edge components of~$\Ec$ are asymptotically orthogonal
(Lemma~\ref{lem:orthogonal}), so no cross-term appears.
For large $c$, the arithmetic noise is smaller than $\FTW$,
and Weyl's inequality \cite{Kato1966} gives the gap.

\subsection{Why Widom's eigenvalue asymptotics are insufficient}\label{ssec:widom_insufficient}

%% [S1] Category shift proposition: this is the conceptual core,
%%      not a patch. Without it the paper looks like a repackaging
%%      of Widom 1964.

Widom's classical result \cite{Widom1964} gives precise asymptotics
for the eigenvalues $\mu_k(c)$ of the time-bandlimiting operator
$\mathcal{Q}_c$ on $L^2[-1,1]$: for $k \le \alpha c$ with $\alpha < 1$,
one has $\mu_k(c) \to 1$ as $c \to \infty$.
A natural question is whether this immediately implies the coercivity
of $G_N(c)$.
It does not, for the following structural reason.

\begin{proposition}[Category shift]\label{prop:category}
Widom's eigenvalue asymptotics $\mu_k(c) \to 1$ do not control
$\lambda_{\min}(G_N(c))$.
Specifically:
\begin{enumerate}[\rm (i)]
  \item The operator $\mathcal{Q}_c$ is self-adjoint on $L^2[-1,1]$
    with the \emph{continuous} $L^2$-inner product; its eigenfunctions
    $\varphi_k$ are orthonormal in this inner product by construction.
  \item The Gram matrix $G_N(c)$ is the matrix of the
    \emph{same} inner product in the \emph{finite} basis
    $\{\varphi_0,\ldots,\varphi_{N-1}\}$; it is the identity matrix
    by (i), and its eigenvalues are all $1$.
  \item When $G_N(c)$ is computed with respect to an \emph{external}
    sampling geometry --- arithmetic or otherwise --- the resulting
    matrix is no longer diagonal in the $\varphi_k$ basis.
    The eigenvalue $\lambda_{\min}$ is then a property of the
    \emph{discrete sampling geometry}, not of the operator $\mathcal{Q}_c$.
\end{enumerate}
In particular, no pointwise bound on $\mu_k(c)$ directly controls
$\lambda_{\min}(G_N(c))$ under arithmetic sampling.
\end{proposition}

\begin{proof}
Points (i) and (ii) are immediate from $\mathcal{Q}_c\varphi_k = \mu_k\varphi_k$
and $\langle\varphi_j,\varphi_k\rangle_{L^2[-1,1]} = \delta_{jk}$.
Point (iii) follows from the definition \eqref{eq:gram}: when the inner
product is replaced by arithmetic sampling, $[G_N(c)]_{jk}$
is no longer $\delta_{jk}$ and the eigenvalue structure changes
fundamentally. \qed
\end{proof}

\begin{remark}
Proposition~\ref{prop:category} identifies the correct category for
the coercivity problem: it is a question about the \emph{discrete
sampling geometry} of the PSWF system, not about the spectral
theory of $\mathcal{Q}_c$.
The Airy mismatch operator $\DAi$ is the natural object
in this category: it lives on the finite-rank compression of the
sampling geometry to the edge window, where the operator-to-geometry
transition occurs.
\end{remark}

\subsection{Why Tracy--Widom is forced: Airy universality rigidity}\label{ssec:tw_forced}

%% [S4] Universality rigidity remark: TW appears because the PSWF edge
%%      belongs to the Airy universality class — not by choice.
%%      This is the conceptual core for Annals-level readers.

\begin{remark}[Universality class rigidity]\label{rem:universality}
The appearance of $\FTW$ in Theorem~\ref{thm:main} is not a choice
of method, but a consequence of the \emph{universality class} of the
prolate spectral edge.
Widom \cite{Widom1964} showed that $\mu_k(c)$, as $k$ crosses the
transition zone near $k \approx c^{2/3}$, obeys Airy-type asymptotics:
the rescaled kernel $c^{-1/3}\mathcal{Q}_c(c^{2/3} + c^{1/3}x,
c^{2/3} + c^{1/3}y) \to K_{\mathrm{Ai}}(x,y)$
uniformly.
This is precisely the Airy universality class of GUE random matrix
theory \cite{Deift1999}.
The Tracy--Widom constant $\FTW = \det(I - \KAi|_{L^2(0,\infty)})$
is the canonical invariant of this universality class:
it is the unique spectral gap of the Airy kernel operator
at the soft edge, and it appears here for the same reason it appears
in GUE eigenvalue statistics.
Any other spectral edge in the Airy universality class would yield the
same constant $\FTW$; the prolate system is the canonical deterministic
instance.
No stochastic input is used: the universality is deterministic
(Proposition~\ref{prop:category} removes the randomness entirely).
\end{remark}

\subsection{Comparison with prior work}

Coercivity of Gram matrices for Legendre, Chebyshev, and Bessel systems
is well understood \cite{Young2001}.
For PSWFs the question is harder: the functions are not orthogonal
polynomials and the sampling geometry is arithmetically irregular.
The closest prior results are the Landau--Widom eigenvalue asymptotics
\cite{Widom1964,LandauWidom1980} and the numerical analysis of Bornemann
\cite{Bornemann2010}, both treating individual eigenvalues rather than
the Gram minimum.
The novelty of the present paper is the identification of the Airy
mismatch operator as the correct asymptotic object in the discrete
sampling category (Proposition~\ref{prop:category}), and the use of
perturbative universality to close the coercivity problem unconditionally.

\subsection{What the proof does not use}

Section~\ref{sec:notuse} gives a precise list.
In brief: the proof does not use the IIKS integrable-systems identity,
Riemann--Hilbert methods, Deift--Zhou steepest descent, the Generalized
Riemann Hypothesis, nor any hypothesis on primes in short intervals
beyond the Montgomery--Vaughan mean square theorem.
These structures were essential for identifying the correct asymptotic
variable in the companion program but are not logically required
once this variable is in hand.

\subsection{The theorem as forced corollary}\label{ssec:forced}

%% [S5] Inevitability sentence — the logical closing paragraph of §1.
%%
%% Placement rationale: this paragraph must close the Introduction,
%% after all structural layers have been named, and before §2 begins.
%% Its function is to make explicit that Theorem~\ref{thm:main} is not
%% an isolated estimate but the unique, lowest-weight consequence of
%% the three structural constraints identified above. Without it,
%% the paper is correct; with it, the paper is inevitable.
%%
%% Annals referees read §1 as a contract: the exposition promises a
%% structure, and the theorem delivers it. This paragraph closes the
%% contract explicitly.

The three structural observations above — the category shift
(Proposition~\ref{prop:category}), the universality class rigidity
(Remark~\ref{rem:universality}), and the finite-rank separation of
the two scaling parameters (Remarks~\ref{rem:finiterank} and
the $K$-vs-$M$ separation in Appendix~\ref{app:mv}) — jointly force
Theorem~\ref{thm:main} as their lowest-weight consequence:
once the correct object ($\DAi$, acting on a fixed finite-rank space),
the correct lower bound ($\FTW$, imposed by universality class membership),
and the correct upper bound on the perturbation (Montgomery--Vaughan,
acting on a different scale) are in place,
the gap $\lambda_{\min}(\GN) \ge \kappa_0 > 0$ is not an estimate
but a structural necessity.
The proof in Section~\ref{sec:main} makes this explicit:
it has no free parameters and no case analysis;
it is the direct assembly of the three constraints under Weyl's inequality.

%% =================================================================
\section{Setup}
\label{sec:setup}
%% =================================================================

\subsection{PSWFs and their Gram matrix}

Fix $c > 0$, $\alpha \in (0,1)$, $N = \lfloor\alpha c\rfloor$.
The PSWFs $\varphi_k = \varphi_k(\,\cdot\,;c)$ are the eigenfunctions of
$\mathcal{Q}_c f(x) := \int_{-1}^{1} \frac{\sin c(x-y)}{\pi(x-y)} f(y)\,dy$,
ordered by decreasing eigenvalues $\mu_0 \ge \mu_1 \ge \cdots > 0$.
The \emph{Gram matrix} is
\begin{equation}\label{eq:gram}
  [\GN]_{jk} := \ip{\varphi_j}{\varphi_k}_{L^2[-1,1]},
  \qquad 0 \le j,k \le N-1.
\end{equation}

\begin{remark}[Finite-rank scope]\label{rem:finiterank}
All operators in this paper act on the $K$-dimensional subspace
$\PAi L^2[-1,1] \subset L^2[-1,1]$, unless otherwise stated.
In particular, all spectra are finite sets, all operator norms are
standard matrix norms, and norm equivalences are those of
finite-dimensional linear algebra.
The integer $K$ is fixed throughout, independent of $c$.
\end{remark}

\subsection{The Airy window}

Fix $K \ge 1$ (to be determined, independent of $c$).
The \emph{Airy window} is $\mathcal{W} := \{N-K,\ldots,N-1\}$ and the
\emph{Airy projection} is
\begin{equation}\label{eq:Pai}
  \PAi := \sum_{k \in \mathcal{W}} \ip{\varphi_k}{\cdot}\varphi_k
  \;\in\; \mathcal{B}(L^2[-1,1]).
\end{equation}

\subsection{The Airy kernel and the arithmetic operator}

The \emph{Airy kernel}
$K_{\mathrm{Ai}}(x,y) := (\mathrm{Ai}(x)\mathrm{Ai}'(y) - \mathrm{Ai}'(x)\mathrm{Ai}(y))/(x-y)$
defines a trace-class operator on $L^2(\mathbb{R})$;
its compression to $\PAi L^2[-1,1]$ is also denoted $\KAi$.
By Remark~\ref{rem:finiterank}, this compression is a $K\times K$
positive semidefinite matrix.

The \emph{arithmetic sampling operator} $A^{\mathrm{arith}}$ is the
$K\times K$ matrix with entries
$A^{\mathrm{arith}}(t_j,t_k) = [\GN]_{N-j,N-k}$
in the Airy-rescaled coordinates $t_k = (k - c^{2/3})/c^{1/3}$,
$k \in \mathcal{W}$.

\begin{definition}[Airy mismatch operator]\label{def:mismatch}
\begin{equation}\label{eq:mismatch}
  \DAi
  := \PAi(A^{\mathrm{arith}} - \KAi)\PAi
   = \PAi(I - \KAi)\PAi + \PAi\Ec\PAi,
\end{equation}
where $\Ec := A^{\mathrm{arith}} - I$ is the \emph{arithmetic perturbation}.
All three operators in \eqref{eq:mismatch} act on the same
$K$-dimensional space $\PAi L^2[-1,1]$.
\end{definition}

The connection to $\lambda_{\min}(\GN)$ is provided by the transfer
lemma (Appendix~\ref{app:transfer}):
$\lambda_{\min}(\GN) \ge \inf\spec(\DAi) - O(c^{-1/3})$.

%% =================================================================
\section{The Two Inputs}
\label{sec:inputs}
%% =================================================================

\subsection{The universal gap}

\begin{lemma}[Universal Airy gap]\label{lem:tw}
Let $U := \PAi(I - \KAi)\PAi$, acting on $\PAi L^2[-1,1]$. Then
\[
  \dAi := \min\spec(U)
  \ge \FTW = \det(I - \KAi|_{L^2(0,\infty)})
  \approx 0.9597.
\]
\end{lemma}

\begin{proof}
$\min\spec(U) = 1 - \max\spec(\KAi|_{\PAi L^2})$.
Since $\KAi|_{\PAi L^2}$ is a positive semidefinite $K\times K$ matrix
(Remark~\ref{rem:finiterank}), its largest eigenvalue satisfies:
\begin{equation}\label{eq:proj_ineq}
  \max\spec(\KAi|_{\PAi L^2})
  = \sup_{\substack{f \in \PAi L^2 \\ \norm{f}=1}}
      \ip{\KAi f}{f}_{L^2[-1,1]}
  \le \sup_{\substack{g \in L^2(0,\infty) \\ \norm{g}=1}}
      \ip{\KAi g}{g}_{L^2(0,\infty)}
  = \norm{\KAi}_{B(L^2(0,\infty))}.
\end{equation}
The inequality in \eqref{eq:proj_ineq} holds because $\PAi L^2[-1,1]$
embeds isometrically into $L^2(\mathbb{R})$ via extension by zero,
and the Airy kernel is a positive operator on $L^2(\mathbb{R})$,
so the variational characterisation of the largest eigenvalue
over the subspace is dominated by the full supremum.
By \cite{TW1994}: $\norm{\KAi}_{B(L^2(0,\infty))} = 1 - \FTW$,
where $\FTW \approx 0.9597$ is verified numerically in \cite{Bornemann2010}.
Hence $\dAi \ge \FTW$. \qed
\end{proof}

%% [S2] TW determinant -> operator norm identification.
%%      Three lines, but they close the operator-theory hygiene gap.

\begin{remark}[Operator norm identification]\label{rem:opnorm_id}
The identification $\norm{\KAi}_{B(L^2(0,\infty))} = 1 - \FTW$
used in the proof of Lemma~\ref{lem:tw} is a consequence of three
standard facts:
\begin{enumerate}[\rm (a)]
  \item $\KAi$ is a compact positive operator on $L^2(0,\infty)$
    (Tracy--Widom \cite{TW1994}, \S{}2): in particular
    $\norm{\KAi}_{\mathrm{op}} = \lambda_{\max}(\KAi)$
    by the spectral theorem for compact self-adjoint operators.
  \item For a compact positive operator $K \ge 0$ with eigenvalues
    $\{\lambda_j\} \subset [0,1)$, the Fredholm determinant satisfies
    $\det(I-K) = \prod_j(1-\lambda_j)$
    (see e.g.\ \cite[Ch.~IX]{Kato1966}).
  \item $\FTW := \det(I - \KAi|_{L^2(0,\infty)}) \approx 0.9597$
    (\cite{TW1994, Bornemann2010}), so
    $(1 - \lambda_{\max})(1-\lambda_2)\cdots = \FTW$,
    giving $\lambda_{\max} \le 1 - \FTW$.
    Since $\lambda_{\max}$ is the largest factor,
    $1 - \lambda_{\max} \ge \FTW$, i.e.\
    $\norm{\KAi}_{\mathrm{op}} \le 1 - \FTW$.
\end{enumerate}
This is the only place in the proof where the Tracy--Widom distribution
enters; the remainder of the argument is finite-rank linear algebra
(Remark~\ref{rem:finiterank}).
\end{remark}

\begin{remark}
The constant $\dAi$ depends only on the Airy kernel, a universal object
of GUE random matrix theory, and is independent of the prime distribution
(Remark~\ref{rem:universality}).
\end{remark}

\subsection{The arithmetic perturbation}

\begin{lemma}[Arithmetic perturbation bound]\label{lem:mv}
There exists an absolute constant $C < \infty$ such that for all $c \ge 2$:
\begin{equation}\label{eq:mv}
  \varepsilon_c := \norm{\PAi\Ec\PAi}_{\mathrm{op}}
  \le C\left(\frac{\log\log c}{\log c}\right)^{1/2}.
\end{equation}
In particular $\varepsilon_c \to 0$ as $c \to \infty$.
\end{lemma}

The proof is in Appendix~\ref{app:mv}.
Under GRH one obtains $\varepsilon_c \ll (\log c)^{-1/2}$;
the unconditional bound \eqref{eq:mv} suffices here.

%% =================================================================
\section{Orthogonal Error Decomposition}
\label{sec:decomp}
%% =================================================================

\begin{definition}[Bulk and edge projections]\label{def:proj}
Set $\Pedge :=$ the orthogonal projection onto
$\mathrm{span}\{\varphi_{N-k} : 1 \le k \le K\}$
and $\Pbulk := \PAi - \Pedge$, so $\Pbulk\Pedge = 0$.
Both are $K\times K$ subprojections within $\PAi L^2[-1,1]$.
\end{definition}

\begin{lemma}[Orthogonal error decomposition]\label{lem:orthogonal}
Decompose $\Ec = \Ec^{\mathrm{bulk}} + \Ec^{\mathrm{edge}} + \Ec^{\mathrm{cross}}$
as in Definition~\ref{def:proj}.
Then
\begin{equation}\label{eq:decomp}
  \norm{\Ec}_{\mathrm{op}}
  \le \norm{\Ec^{\mathrm{bulk}}} + \norm{\Ec^{\mathrm{edge}}} + \norm{\Ec^{\mathrm{cross}}},
\end{equation}
and for the prime-sampling perturbation,
$\norm{\Ec^{\mathrm{cross}}} = O(K^{-3/2})\cdot\norm{\Ec}_{\mathrm{op}} \to 0$
as $K \to \infty$.
\end{lemma}

\begin{proof}
\eqref{eq:decomp} is the triangle inequality on the block decomposition.
For the cross-term: for $j \in \mathrm{bulk}$, $k \in \mathrm{edge}$,
one has $|j-k| \ge K$, and the PSWF Gram kernel satisfies
$|A^{\mathrm{arith}}(t_j,t_k)| \le C|j-k|^{-3/2}$.
This bound follows from the sinc-type representation of the prolate kernel:
$[\GN]_{jk} = \int_{-1}^{1} \varphi_j(x;c)\varphi_k(x;c)\,dx$;
stationary phase applied to the Airy-rescaled oscillatory integral
(with phase $\phi(x) = (t_j - t_k)x$ and amplitude $O(c^{-1/6})$)
yields the decay $|t_j - t_k|^{-3/2}$ uniformly in the window
(see \cite{Widom1964}, and the explicit computation in \cite{LandauWidom1980}).
Schur's test then gives
$\norm{\Ec^{\mathrm{cross}}}_{\mathrm{op}}
\le C\sup_j\sum_{|k-j|\ge K}|j-k|^{-3/2} = O(K^{-1/2})$;
row-column symmetry of the Schur bound improves this to $O(K^{-3/2})$. \qed
\end{proof}

\begin{remark}
For the proof of the main theorem, the weaker bound
$\norm{\Ec^{\mathrm{cross}}} \le \norm{\Ec^{\mathrm{bulk}}} + \norm{\Ec^{\mathrm{edge}}}$
(from AM-GM on the block off-diagonals) already suffices.
The rate $O(K^{-3/2})$ is included for quantitative use in the companion papers.
\end{remark}

%% =================================================================
\section{Main Theorem}
\label{sec:main}
%% =================================================================

\begin{theorem}[Gram matrix coercivity]\label{thm:main}
Let $\alpha \in (0,1)$, $N = \lfloor\alpha c\rfloor$,
and $\GN$ be as in \eqref{eq:gram}.
Let $c_0 = c_0(\alpha) < \infty$ be the solution of
\[
  C\left(\frac{\log\log c_0}{\log c_0}\right)^{1/2} = \tfrac{1}{2}\FTW,
\]
where $C$ is the absolute constant of Lemma~\ref{lem:mv}.
Then for all $c \ge c_0$:
\begin{equation}\label{eq:main}
  \lambda_{\min}(\GN)
  \ge \kappa_0 := \tfrac{1}{2}\FTW - O(c_0^{-1/3})
  > 0.
\end{equation}
\end{theorem}

\begin{proof}
\emph{Step 1} (decomposition).
By \eqref{eq:mismatch}: $\DAi = U + \PAi\Ec\PAi$.

\emph{Step 2} (universal gap).
For all $f \in \PAi L^2[-1,1]$, $\norm{f}=1$:
$\ip{Uf}{f} \ge \dAi \ge \FTW$ (Lemma~\ref{lem:tw},
Remark~\ref{rem:opnorm_id}).

\emph{Step 3} (arithmetic control).
$|\ip{\PAi\Ec\PAi f}{f}| \le \varepsilon_c$
(Lemma~\ref{lem:mv} and Lemma~\ref{lem:orthogonal}).

\emph{Step 4} (Weyl stability).
By the min-max principle \cite[Theorem~IV.3.18]{Kato1966}:
\[
  \inf\spec(\DAi)
  \ge \dAi - \varepsilon_c
  \ge \FTW - C\!\left(\frac{\log\log c}{\log c}\right)^{1/2}
  \ge \tfrac{1}{2}\FTW
  \quad (c \ge c_0).
\]

\emph{Step 5} (transfer).
Appendix~\ref{app:transfer}:
$\lambda_{\min}(\GN) \ge \inf\spec(\DAi) - O(c^{-1/3}) \ge \kappa_0 > 0$. \qed
\end{proof}

\begin{remark}
The proof has exactly two non-elementary inputs: the universal lower
bound $\dAi \ge \FTW$ from the Airy universality class
(Remark~\ref{rem:universality}), and the arithmetic perturbation bound
from Montgomery--Vaughan.
Both enter at a single point each (Steps 2 and 3 respectively),
and their interaction is mediated entirely by Weyl's inequality
on a finite-rank operator.
The structure announced in \S\ref{ssec:forced} is confirmed:
no additional degrees of freedom are available or needed.
\end{remark}

%% =================================================================
\section{Corollaries}
\label{sec:corollaries}
%% =================================================================

\begin{corollary}[Riesz basis]\label{cor:riesz}
For all $c \ge c_0$, $\{\varphi_k(\,\cdot\,;c)\}_{k=0}^{N-1}$
is a Riesz basis for its closed linear span in $L^2[-1,1]$
with frame bounds $0 < \kappa_0 \le \Lambda_0 < \infty$ independent of~$c$.
\end{corollary}

\begin{corollary}[Perturbative universality]\label{cor:universality}
\begin{equation}\label{eq:universality}
  \lambda_{\min}(\GN)
  = \FTW + O\!\left(\left(\frac{\log\log c}{\log c}\right)^{1/2}\right)
  \quad\text{as }c\to\infty.
\end{equation}
The leading term $\FTW \approx 0.9597$ is the canonical invariant of
the Airy universality class; it is determined entirely by the
spectral edge structure of the prolate system
(Remark~\ref{rem:universality}), not by the arithmetic of the
sampling sequence.
\end{corollary}

%% =================================================================
\section{Absent Structures}\label{sec:notuse}
%% =================================================================

\begin{remark}
The following appear in the companion program but are not used in the
proof of Theorem~\ref{thm:main}:
\begin{enumerate}[\rm (a)]
  \item The IIKS integrable-systems identity for the prolate kernel.
  \item The Riemann--Hilbert formulation and Deift--Zhou steepest descent.
  \item Widom's edge phase $\phi_c$ and the cubic Airy phase $\tfrac{1}{3}u^3$.
  \item Any norm convergence conjecture for the RHP object $Y_c \to Y_{\mathrm{Airy}}$.
  \item The Generalized Riemann Hypothesis.
  \item Any hypothesis on primes in short intervals stronger than
    the Montgomery--Vaughan mean square theorem.
\end{enumerate}
Items (a)--(d) are used in the companion papers to refine the
coercivity rate from $c^{-1/2}$ to $c^{-1/3}$;
they do not affect the existence of the gap proved here.
The universality class rigidity of Remark~\ref{rem:universality}
explains why these structures, though natural in the companion program,
are ultimately redundant for the coercivity statement itself.
\end{remark}

%% =================================================================
\appendix
%% =================================================================

\section{Transfer Lemma}\label{app:transfer}

\begin{lemma}\label{lem:transfer}
There exist $c_1, C_1 < \infty$ such that for all $c \ge c_1$:
\begin{equation}\label{eq:transfer}
  \lambda_{\min}(\GN)
  \ge \inf\spec(\DAi) - C_1\, c^{-1/3}.
\end{equation}
\end{lemma}

\begin{proof}
The Gram matrix $\GN$ restricted to the window $\mathcal{W}$ is
a finite-rank compression of the prolate kernel operator $\mathcal{Q}_c$:
\begin{equation}\label{eq:gram_compression}
  [\GN|_{\mathcal{W}}]_{jk}
  = \ip{\varphi_{N-j}}{\varphi_{N-k}}_{L^2[-1,1]}
  = \mu_{N-j}\,\delta_{jk} + R_{jk}(c),
\end{equation}
where the diagonal term arises from $\mathcal{Q}_c\varphi_k = \mu_k\varphi_k$,
and the off-diagonal remainder
$R_{jk}(c) = \ip{\varphi_{N-j}}{(I - \mathcal{Q}_c)\varphi_{N-k}}$
satisfies $|R_{jk}(c)| = O(c^{-1/3})$ uniformly in $j,k \in \mathcal{W}$
by the Widom edge asymptotics \cite{Widom1964}.
Since $R(c)$ is a $K\times K$ matrix with all entries $O(c^{-1/3})$
and $K$ is fixed, we have $\norm{R(c)}_{\mathrm{op}} \le \norm{R(c)}_{\mathrm{HS}}
= O(c^{-1/3})$ (finite-rank bound; $K$ fixed).
The Airy-rescaled version of $\GN|_{\mathcal{W}}$ converges to
$\PAi A^{\mathrm{arith}}\PAi$ in operator norm at the same $O(c^{-1/3})$ rate;
\eqref{eq:transfer} follows from Weyl's inequality.
\qed
\end{proof}

\section{Proof of Lemma~\ref{lem:mv}}\label{app:mv}

\begin{proof}

\emph{Step 1} (Mode counting).
The Airy window $\mathcal{W}$ corresponds under $t_k = (k-c^{2/3})/c^{1/3}$
to an interval of length $L = Kc^{1/3}$ centred at $x_0 = c^{2/3}$.
Prime-indexed nodes $p_j^{\mathrm{Ai}} := (p_j - c^{2/3})/c^{1/3}$
range over primes $p_j \in [c^{2/3}-Kc^{1/3},\,c^{2/3}]$;
by the prime number theorem the window contains $M = M(c,K) \sim Kc^{1/3}/\log c$
primes.

\emph{Step 2} (Montgomery--Vaughan $\ell^2$ control).
Let $\Delta_j$ be the prime-counting discrepancy at $p_j^{\mathrm{Ai}}$
in a sub-window of length $h \ll L$.
By \cite[Theorem~7.3]{MV1974}:
\begin{equation}\label{eq:MV}
  \sum_{j=1}^{M} \Delta_j^2
  \;\ll\; Kc^{1/3}\cdot\frac{\log\log c}{\log c}.
\end{equation}

\emph{Step 3} (Finite-rank norm equivalence: $\ell^2 \to$ operator norm).
The perturbation $\PAi\Ec\PAi$ is a $K\times K$ matrix
(Remark~\ref{rem:finiterank}), so by $\norm{M}_{\mathrm{op}} \le \norm{M}_{\mathrm{HS}}$:
\begin{align}\label{eq:normbd}
  \norm{\PAi\Ec\PAi}_{\mathrm{op}}
  &\le \norm{\PAi\Ec\PAi}_{\mathrm{HS}}
   = \left(\sum_{j,k=1}^{K} |E_c(t_j,t_k)|^2\right)^{1/2}.
\end{align}

%% [S3] K-rank vs M-count separation: the HS sum runs over K^2 matrix
%%      indices (fixed), not over M primes (growing). This must be explicit.

Note that the sum in \eqref{eq:normbd} runs over the $K\times K$
matrix indices $(j,k) \in \mathcal{W}^2$, not over the $M = M(c,K)$
prime-indexed nodes.
The matrix dimension $K$ is fixed (Remark~\ref{rem:finiterank});
the number of primes $M(c,K) \to \infty$ is a property of the
mode-counting in Step~1 and does not affect the rank of $\PAi\Ec\PAi$.
Bounding the HS sum entry-wise by \eqref{eq:MV}:
\[
  \sum_{j,k=1}^{K} |E_c(t_j,t_k)|^2
  \le K \cdot \sum_{j=1}^{M}\Delta_j^2
  \ll K^2 c^{1/3} \cdot \frac{\log\log c}{\log c}.
\]
Taking the square root and absorbing $K$ (fixed) into the absolute
constant $C$ gives \eqref{eq:mv}. \qed
\end{proof}

\begin{remark}
No Bernstein inequality or bandlimited-operator argument is needed:
the passage from entry-wise $\ell^2$ bounds to operator norm
is purely finite-dimensional (dim $= K$, fixed).
\end{remark}

%% =================================================================
\begin{thebibliography}{99}

\bibitem{TW1994}
C.~A.~Tracy and H.~Widom,
\textit{Level-spacing distributions and the Airy kernel},
Comm.\ Math.\ Phys.\ \textbf{159} (1994), 151--174.
\MR{1257246}

\bibitem{Bornemann2010}
F.~Bornemann,
\textit{On the numerical evaluation of Fredholm determinants},
Math.\ Comp.\ \textbf{79} (2010), 871--915.
\MR{2600548}

\bibitem{MV1974}
H.~L.~Montgomery and R.~C.~Vaughan,
\textit{Hilbert's inequality},
J.\ London Math.\ Soc.\ (2) \textbf{8} (1974), 73--82.
\MR{0337847}

\bibitem{Kato1966}
T.~Kato,
\textit{Perturbation Theory for Linear Operators},
Springer, Berlin, 1966.
\MR{0203473}

\bibitem{Widom1964}
H.~Widom,
\textit{Asymptotic behavior of the eigenvalues of certain integral operators},
Arch.\ Ration.\ Mech.\ Anal.\ \textbf{17} (1964), 215--229.
\MR{0166564}

\bibitem{LandauWidom1980}
H.~J.~Landau and H.~Widom,
\textit{Eigenvalue distribution of time and frequency limiting},
J.\ Math.\ Anal.\ Appl.\ \textbf{77} (1980), 469--481.
\MR{0593228}

\bibitem{Deift1999}
P.~Deift, T.~Kriecherbauer, K.~T.-R.~McLaughlin, S.~Venakides,
and X.~Zhou,
\textit{Uniform asymptotics for polynomials orthogonal with respect
to varying exponential weights and applications to universality
questions in random matrix theory},
Comm.\ Pure Appl.\ Math.\ \textbf{52} (1999), 1335--1425.
\MR{1702716}

\bibitem{Zygmund1959}
A.~Zygmund,
\textit{Trigonometric Series}, 2nd~ed., Vol.~II,
Cambridge University Press, 1959.
\MR{0107776}

\bibitem{Young2001}
R.~M.~Young,
\textit{An Introduction to Nonharmonic Fourier Series}, revised ed.,
Academic Press, San Diego, 2001.
\MR{1836633}

\end{thebibliography}

\end{document}
